% !TeX root = ../main.tex
% 按课件顺序整理；原有五处 TODO 及六条补写注释已落实。
% 页码默认指所提供 PDF 的文件页序，避免与教材印刷页码混淆。
\section{计算流体力学基础}
\label{sec:basics}

本讲沿着“控制方程 $\to$ PDE 性质 $\to$ 离散与稳定性 $\to$ 计算网格”展开。
CFD 用离散方程近似连续流动；计算结果需要解析解、网格加密研究和实验数据的检验。
DNS 直接解析相关湍流尺度，LES 解析大尺度并模拟亚格子作用，RANS 求解平均场并封闭雷诺应力。
数据驱动模型同样依赖数据质量与物理检验。\cite{cfd:ppt}

\subsection{控制方程和矢量形式控制方程}
\label{sec:cfd-equations}
\noindent\textit{依据：课件第18--26页；Anderson 第2.3--2.10节。}

\paragraph{记号与假设}
考虑连续介质，速度 $\mathbf v=(u,v,w)^{\mathsf T}$，密度 $\rho$，压强 $p$，
温度 $T$，比内能 $e$，单位质量体力 $\mathbf g$。物质导数为
\[
  \frac{D}{Dt}=\frac{\partial}{\partial t}+\mathbf v\cdot\nabla,
  \qquad E=e+\frac12|\mathbf v|^2,\qquad E_t=\rho E.
\]
$E$ 是单位质量总能，$E_t$ 是单位体积总能，均不含体力势能。
$\dot q$ 是单位质量、单位时间的外部加热率；$\mathbf q$ 是热通量，二者不同。
单组分 Fourier 导热给出 $\mathbf q=-k\nabla T$，$k$ 为导热系数。
对牛顿流体，黏性应力为
\begin{equation}
 \boldsymbol\tau=\mu\bigl[\nabla\mathbf v+(\nabla\mathbf v)^{\mathsf T}\bigr]
       +\lambda_b(\nabla\cdot\mathbf v)\mathbf I.
 \label{eq:cfd-stress}
\end{equation}
这里 $\tau_{ij}$ 表示法向为 $j$ 的面上沿 $i$ 方向的黏性应力，且
$\tau_{ij}=\tau_{ji}$；$\lambda_b$ 为第二黏性系数。
Stokes 假设取 $\lambda_b=-2\mu/3$，这不是所有流体均成立的恒等式。

\paragraph{A. 非守恒形式：跟随流体微团}
质量、动量、内能和组分输运分别满足
\begin{align}
 &\frac{D\rho}{Dt}+\rho\nabla\cdot\mathbf v=0,
 \label{eq:cfd-mass-material}\\
 &\rho\frac{D\mathbf v}{Dt}=\rho\mathbf g-\nabla p+\nabla\cdot\boldsymbol\tau,
 \label{eq:cfd-momentum-material}\\
 &\underbrace{\rho\frac{De}{Dt}}_{\text{内能变化}}
   =\underbrace{\rho\dot q}_{\text{体热源}}
    \underbrace{-p\nabla\cdot\mathbf v}_{\text{压缩功}}
    \underbrace{-\nabla\cdot\mathbf q}_{\text{热通量净输入}}
    +\underbrace{\boldsymbol\tau:\nabla\mathbf v}_{\text{黏性耗散}},
 \label{eq:cfd-internal-energy}\\
 &\rho\frac{DY_s}{Dt}=-\nabla\cdot\mathbf j_s+\dot\omega_s,
       \qquad s=1,\ldots,N_s.
 \label{eq:cfd-species-material}
\end{align}
\begin{itemize}
 \item 连续方程：$D\rho/Dt$ 为微团密度变化，$\rho\nabla\cdot\mathbf v$ 为体积膨胀的贡献。
       不可压条件为 $D\rho/Dt=0$，结合连续方程得到 $\nabla\cdot\mathbf v=0$。
 \item 动量方程：左端为惯性项；右端依次是体力、压力梯度力和黏性力，均按单位体积计。
 \item 内能方程：压缩时 $\nabla\cdot\mathbf v<0$，压力功增加内能；
       单组分时 $-\nabla\cdot\mathbf q=\nabla\cdot(k\nabla T)$。
       $\boldsymbol\tau:\nabla\mathbf v=\sum_{i,j}\tau_{ij}\partial v_i/\partial x_j$
       是双点积，不能写成含义不明的“$\boldsymbol\tau\cdot\nabla\cdot\mathbf v$”。
 \item 组分方程：$Y_s$ 为质量分数，$\mathbf j_s=\rho Y_s\mathbf V_s$ 为相对质量平均速度的扩散通量，
       $\mathbf V_s$ 为扩散速度；$\dot\omega_s$ 是单位体积化学反应质量源。
       满足 $\sum_sY_s=1$、$\sum_s\mathbf j_s=0$、$\sum_s\dot\omega_s=0$，故只有 $N_s-1$ 个独立组分方程。
\end{itemize}
多组分时，忽略其他交叉输运效应，可取
$\mathbf q=-k\nabla T+\sum_s h_s\mathbf j_s$，其中 $h_s$ 为组分比焓。
若 $e$ 已包含生成内能，化学能转化通过组分变化体现，不另加重复的反应热源；
若使用显内能，则须配套写出化学热源。\cite{cfd:ppt,cfd:anderson}

\paragraph{B. 守恒形式：局部积累与净流出}
对任意足够光滑的单位质量变量 $\phi$，由连续方程有
\begin{equation}
 \frac{\partial(\rho\phi)}{\partial t}
 +\nabla\cdot(\rho\mathbf v\phi)=\rho\frac{D\phi}{Dt}.
 \label{eq:cfd-conservative-identity}
\end{equation}
因此可以逐项改写为
\begin{align}
 &\rho_t+\nabla\cdot(\rho\mathbf v)=0,\\
 &(\rho\mathbf v)_t+\nabla\cdot(\rho\mathbf v\otimes\mathbf v+p\mathbf I-\boldsymbol\tau)
       =\rho\mathbf g,\\
 &(\rho e)_t+\nabla\cdot(\rho e\mathbf v)
       =\rho\dot q-p\nabla\cdot\mathbf v-\nabla\cdot\mathbf q
                       +\boldsymbol\tau:\nabla\mathbf v,\\
 &(\rho Y_s)_t+\nabla\cdot(\rho Y_s\mathbf v+\mathbf j_s)=\dot\omega_s.
\end{align}
这里 $\otimes$ 表示张量积。内能的输运虽已写成散度形式，压力功仍不在通量中；
CFD 的\textbf{强守恒能量方程通常使用总能}。将动量方程与速度点乘得到动能方程，再与内能方程相加，有
\begin{equation}
 \frac{\partial E_t}{\partial t}
 +\nabla\cdot\bigl[(E_t+p)\mathbf v-\boldsymbol\tau\cdot\mathbf v+\mathbf q\bigr]
 =\rho\mathbf g\cdot\mathbf v+\rho\dot q.
 \label{eq:cfd-total-energy}
\end{equation}
其结构是
\[
 \boxed{\text{单位体积守恒量的局部变化}+\text{通量的净流出}=\text{源项}}.
\]
例如，固定控制体 $\Omega$ 上的质量守恒为
\[
 \frac{\dd}{\dd t}\int_\Omega\rho\,\dd V
 +\int_{\partial\Omega}\rho\mathbf v\cdot\mathbf n\,\dd A=0.
\]
固定控制体与随体微团是两种推导出发点；\textbf{守恒性取决于方程结构，不能仅由坐标是否移动判定}。
光滑解下两种微分形式等价；有激波时，守恒律应按积分／弱形式理解，
不能未经论证地用非守恒乘积替代，否则可能给出错误的激波速度。

\paragraph{C. 矢量守恒形式与通量分解}
先考虑单组分三维流动。取
\begin{equation}
 \mathbf U=\begin{bmatrix}\rho&\rho u&\rho v&\rho w&E_t\end{bmatrix}^{\mathsf T},
 \qquad
 \mathbf U_t+\mathbf F_x+\mathbf G_y+\mathbf H_z=\mathbf S.
 \label{eq:cfd-vector-conservation}
\end{equation}
$\mathbf U$ 为守恒变量（conservative variables）；
$\mathbf S$ 是\textbf{所有体源项的向量}，并非专指组分变化。
本讲用 $\mathbf S$ 代替课件中的源项符号 $J$，以免与网格 Jacobian 混淆。
在上述假设下，
\[
 \mathbf S=\begin{bmatrix}
 0&\rho g_x&\rho g_y&\rho g_z&\rho\mathbf g\cdot\mathbf v+\rho\dot q
 \end{bmatrix}^{\mathsf T}.
\]
按课件约定，$\mathbf F=\mathbf F_I-\mathbf F_V$、
$\mathbf G=\mathbf G_I-\mathbf G_V$、$\mathbf H=\mathbf H_I-\mathbf H_V$。
其中 $I$ 表示无黏部分，$V$ 表示黏性与热输运部分：
\begin{equation}
 \mathbf F_I=\begin{bmatrix}\rho u\\\rho u^2+p\\\rho uv\\\rho uw\\u(E_t+p)\end{bmatrix},\quad
 \mathbf G_I=\begin{bmatrix}\rho v\\\rho uv\\\rho v^2+p\\\rho vw\\v(E_t+p)\end{bmatrix},\quad
 \mathbf H_I=\begin{bmatrix}\rho w\\\rho uw\\\rho vw\\\rho w^2+p\\w(E_t+p)\end{bmatrix}.
 \label{eq:cfd-inviscid-fluxes}
\end{equation}
\begin{equation}
 \mathbf F_V=\begin{bmatrix}0\\\tau_{xx}\\\tau_{yx}\\\tau_{zx}\\\Psi_x\end{bmatrix},\quad
 \mathbf G_V=\begin{bmatrix}0\\\tau_{xy}\\\tau_{yy}\\\tau_{zy}\\\Psi_y\end{bmatrix},\quad
 \mathbf H_V=\begin{bmatrix}0\\\tau_{xz}\\\tau_{yz}\\\tau_{zz}\\\Psi_z\end{bmatrix},
 \label{eq:cfd-viscous-fluxes}
\end{equation}
其中能量输运分量全部展开为
\begin{align*}
 \Psi_x&=u\tau_{xx}+v\tau_{yx}+w\tau_{zx}-q_x,\\
 \Psi_y&=u\tau_{xy}+v\tau_{yy}+w\tau_{zy}-q_y,\\
 \Psi_z&=u\tau_{xz}+v\tau_{yz}+w\tau_{zz}-q_z.
\end{align*}
由于总通量是 $I-V$，最终能量通量中的导热项为 $+q_x=-kT_x$，其余方向同理。
忽略黏性、热传导及体源即得到 Euler 方程。量热完全气体可用
$p=(\gamma-1)[E_t-\rho(u^2+v^2+w^2)/2]$ 闭合，$\gamma=c_p/c_v$ 为常数。
增加组分时，把 $\rho Y_s$ 加入 $\mathbf U$，相应 $x$ 向总通量为
$\rho uY_s+j_{s,x}$，故按同一减号约定，其 $V$ 分量是 $-j_{s,x}$，源项是 $\dot\omega_s$。

\paragraph{初始与边界条件}
同一组控制方程可描述不同流动，差别来自几何、初值与边界条件。
黏性固壁通常施加无滑移和不可穿透条件；Euler 固壁只施加不可穿透条件。
热边界可取给定壁温或给定热流，绝热对应 $\mathbf q\cdot\mathbf n=0$。
给定壁温并非所有固壁的统一条件。

\subsection{偏微分方程的数学性质}
\label{sec:cfd-pde}
\noindent\textit{依据：课件第27--35页；Anderson 第3.2--3.5节。}

本节分析信息如何传播，以及怎样配置初始与边界数据。
拟线性指\textbf{最高阶导数线性出现}，其系数可以依赖自变量、未知量及低阶导数。
以下二元一阶方程组是讨论特征的模型：
\begin{equation}
 \begin{aligned}
 a_1u_x+b_1u_y+c_1v_x+d_1v_y&=f_1,\\
 a_2u_x+b_2u_y+c_2v_x+d_2v_y&=f_2.
 \end{aligned}
 \label{CFD:LinearEquationGroup}
\end{equation}
这里 $u,v$ 是一般未知函数。Euler 方程可写成一阶拟线性系统；
Navier--Stokes 方程包含黏性二阶导数，不能直接全部归为这个一阶模型。

\paragraph{特征线、依赖域与影响域}
沿非特征曲线给定数据时，方程通常能确定横向导数；沿特征曲线，相应系数矩阵退化，
数据必须满足相容关系。对 $u_t+cu_x=0$，
\[
 \frac{\dd x}{\dd t}=c,\qquad
 \frac{\dd u}{\dd t}=u_t+cu_x=0,
 \qquad u(x,t)=u_0(x-ct).
\]
即信息沿 $x-ct=\text{常数}$ 传播。只有把某个自变量选为时间后，特征斜率才具有传播速度的解释。
某点的\textbf{依赖域}是决定其解值所需的过去数据区域；\textbf{影响域}是该点扰动可到达的未来区域。

\begin{figure}[htbp]
 \centering
 \setlength{\unitlength}{1cm}
 \begin{picture}(10,4.4)
  \put(0.5,0.7){\vector(1,0){8.7}}
  \put(0.5,0.7){\vector(0,1){3.3}}
  \put(9.3,0.65){$x$}\put(0.25,4){$t$}
  \put(4.7,3.5){\circle*{0.10}}
  \put(4.85,3.5){$P=(x_0,t_0)$}
  \put(1.9,0.7){\line(1,1){2.8}}
  \put(4.7,3.5){\line(1,-1){2.8}}
  \linethickness{1pt}
  \put(1.9,0.7){\line(1,0){5.6}}
  \put(1.4,0.15){\small $x_0-ct_0$}
  \put(7.0,0.15){\small $x_0+ct_0$}
  \put(2.0,2.5){\small $\dd x/\dd t=c$}
  \put(6.1,2.5){\small $\dd x/\dd t=-c$}
  \put(3.6,1.25){\small 过去依赖域}
 \end{picture}
 \caption{一维波动方程 $u_{tt}=c^2u_{xx}$（$c>0$）的特征线与依赖域。
 初始线上的粗线段给出决定 $P$ 所需的数据范围；交换过去与未来即得到影响域。
 单向对流方程在初始线上只依赖一个点，不能与双向波动方程混同。}
 \label{fig:cfd-characteristics}
\end{figure}

\paragraph{A. 判别式法}
在曲线上加入 $\dd u=u_x\dd x+u_y\dd y$、$\dd v=v_x\dd x+v_y\dd y$，
与式~\eqref{CFD:LinearEquationGroup} 联立为
\begin{equation}
 \underbrace{\begin{bmatrix}
 a_1&b_1&c_1&d_1\\a_2&b_2&c_2&d_2\\
 \dd x&\dd y&0&0\\0&0&\dd x&\dd y
 \end{bmatrix}}_{\mathcal M}
 \underbrace{\begin{bmatrix}u_x\\u_y\\v_x\\v_y\end{bmatrix}}_{\mathbf z}
 =\underbrace{\begin{bmatrix}f_1\\f_2\\\dd u\\\dd v\end{bmatrix}}_{\mathbf b}.
 \label{eq:cfd-characteristic-matrix}
\end{equation}
当 $\det\mathcal M\ne0$ 时，四个偏导数可被唯一恢复；
$\det\mathcal M=0$ 给出特征方向，此时右端还须满足相容条件。
若 $\dd x\ne0$，令 $m=\dd y/\dd x$，得到
\begin{equation}
 am^2+bm+c_0=0,
 \quad\begin{cases}
 a=a_1c_2-a_2c_1,\\
 b=-(a_1d_2-a_2d_1+b_1c_2-b_2c_1),\\
 c_0=b_1d_2-b_2d_1.
 \end{cases}
 \label{eq:cfd-characteristic-polynomial}
\end{equation}
此处用 $c_0$ 区分多项式常数项与波速 $c$。
应保留齐次形式 $a(\dd y)^2+b\dd x\dd y+c_0(\dd x)^2=0$，
以免漏掉 $\dd x=0$ 的竖直特征方向。不恒为零的二次特征式的判别式为 $\Delta=b^2-4ac_0$：
\begin{center}
 \begin{tabular}{lll}
  \toprule
  判别式 & 每点的实特征方向数 & 分类含义\\
  \midrule
  $\Delta>0$ & 两个不同方向 & 双曲型\\
  $\Delta=0$ & 一个重合方向 & 重特征／退化情形，须进一步判断\\
  $\Delta<0$ & 无实特征方向 & 椭圆型\\
  \bottomrule
 \end{tabular}
\end{center}
“方向”沿流场积分后才形成特征曲线，并非整个计算域只存在一条或两条曲线。

\begin{remark}[重特征不等于一般系统必为抛物型]
课件把 $\Delta=0$ 简记为“抛物型”。严格地说，二阶标量方程
$A_2\phi_{xx}+2B_2\phi_{xy}+C_2\phi_{yy}+\text{低阶项}=0$
在主部矩阵非零时，可按 $B_2^2-A_2C_2$ 的正、零、负分别分为双曲、抛物、椭圆型。
对一阶系统，重根本身不足以判定抛物型，还必须检查特征向量与完整微分结构。
若特征多项式恒为零，则上述三分判据不适用。
\end{remark}

\paragraph{B. 特征值法}
定义
\[
 \mathbf W=\begin{bmatrix}u\\v\end{bmatrix},\quad
 A=\begin{bmatrix}a_1&c_1\\a_2&c_2\end{bmatrix},\quad
 B=\begin{bmatrix}b_1&d_1\\b_2&d_2\end{bmatrix},\quad
 \mathbf f=\begin{bmatrix}f_1\\f_2\end{bmatrix}.
\]
原式为 $A\mathbf W_x+B\mathbf W_y=\mathbf f$；
非零源项不能无故丢掉，但不改变这里的主部特征式。
若 $A$ 可逆，令 $C=A^{-1}B$，则
\begin{equation}
 \mathbf W_x+C\mathbf W_y=A^{-1}\mathbf f,
 \qquad \det(B-mA)=\det(A)\det(C-mI)=0.
 \label{eq:cfd-eigen-characteristics}
\end{equation}
因此 $C$ 的特征值就是 $\dd y/\dd x$。
两个不同实特征值对应严格双曲型；一对非实共轭特征值对应这个实二阶系统的椭圆型。
\textbf{负实特征值仍是实特征值，只表示反向特征斜率，绝非椭圆性的判据。}
重实特征值且具有完整特征向量时，仍可为非严格双曲型；例如
$C=cI$ 描述两份独立对流。若特征向量不足，则须分析弱双曲性与适定性。
$A$ 奇异时不可使用 $A^{-1}$，应回到 $\det(B\dd x-A\dd y)=0$。

对多维演化系统 $\mathbf U_t+A\mathbf U_x+B\mathbf U_y=0$，
应对\textbf{每个实方向} $\mathbf n$ 检查 $A_n=n_xA+n_yB$：实谱且存在一致有界的对角化
是强双曲性的条件，仅分别查看 $A$、$B$ 的特征值并不充分。
在 $\rho>0,p>0$ 的理想气体物理状态下，Euler 方程沿单位法向的特征速度为
$v_n-a_s,v_n,v_n,v_n,v_n+a_s$（三维），$a_s$ 为声速；重合的对流特征仍具有完整特征向量。

\begin{example}[课件中的二维定常小扰动流]
可压缩、无黏、无旋的均匀来流小扰动方程为
\[
 (1-M_\infty^2)u'_x+v'_y=0,\qquad u'_y-v'_x=0.
\]
当 $M_\infty\ne1$ 时，$m^2=1/(M_\infty^2-1)$。
故 $M_\infty>1$ 时有两族实特征线，为双曲型；$M_\infty<1$ 时为椭圆型。
在 $M_\infty=1$，矩阵 $A$ 奇异；相应势方程的二阶主部退化，不能继续套用 $A^{-1}B$。
这些结论针对\emph{定常势流小扰动模型}，不能据此称非定常亚声速 Euler 系统为椭圆型。
\end{example}

\paragraph{分类对求解的影响}
双曲型有有限传播速度，需要沿入射特征配置边界数据；
抛物型如热方程 $u_t=\alpha u_{xx}$（$\alpha>0$）随时间扩散，在连续模型中具有无限传播速度；
椭圆型如 Laplace 方程描述整体边值平衡，局部变化通常影响全域。
“上、下游相互影响”应结合特征速度与具体模型判断，不能作为全部 PDE 的共同性质。
适定问题要求解存在、唯一并连续依赖数据；离散格式不能弥补错误的定解条件。

\subsection{基础数值算法}
\label{sec:cfd-numerics}
有限差分直接近似微分算子；有限体积对控制体积分并用共享面通量保证离散守恒；
有限元通常从弱形式与基函数近似出发。以下按课件顺序，以均匀网格上的有限差分为主。

\subsubsection{有限差分方法基础}
\noindent\textit{依据：课件第35--48页；Anderson 第4.2--4.3节。}

记 $x_i=i\Delta x$、$t_n=n\Delta t$，$u_i^n\approx u(x_i,t_n)$。
在函数足够光滑时，Taylor 展开给出
\[
 u(x_i\pm\Delta x)=u_i\pm\Delta x\,u_x
 +\frac{\Delta x^2}{2}u_{xx}\pm\frac{\Delta x^3}{6}u_{xxx}
 +\frac{\Delta x^4}{24}u_{xxxx}+\cdots.
\]
相减消去偶数阶项，相加消去奇数阶项，得到
\begin{align}
 (u_x)_i&=\frac{u_{i+1}-u_i}{\Delta x}+O(\Delta x)
                       &&\text{向前差分},\\
 (u_x)_i&=\frac{u_i-u_{i-1}}{\Delta x}+O(\Delta x)
                       &&\text{向后差分},\\
 (u_x)_i&=\frac{u_{i+1}-u_{i-1}}{2\Delta x}+O(\Delta x^2)
                       &&\text{中心差分},\\
 (u_{xx})_i&=\frac{u_{i+1}-2u_i+u_{i-1}}{\Delta x^2}+O(\Delta x^2)
                       &&\text{二阶导数的中心差分}.
 \label{eq:cfd-second-difference}
\end{align}
边界附近也可使用二阶单边公式
\[
 (u_x)_i=\frac{-3u_i+4u_{i+1}-u_{i+2}}{2\Delta x}+O(\Delta x^2).
\]
二维混合导数可先后应用两个方向的中心差分：
\begin{equation}
 (u_{xy})_{i,j}=\frac{u_{i+1,j+1}-u_{i+1,j-1}-u_{i-1,j+1}+u_{i-1,j-1}}
 {4\Delta x\Delta y}+O(\Delta x^2+\Delta y^2).
 \label{eq:cfd-mixed-difference}
\end{equation}
以上精度针对均匀网格与光滑解；非均匀网格必须重新计算系数，间断处不能直接使用 Taylor 精度结论。
例如前向时间差分的一阶导数误差为 $O(\Delta t)$，不等于每步解的误差也为 $O(\Delta t)$；
把微分方程残差乘以 $\Delta t$ 后，每步局部误差还会多一个时间步因子。

\subsubsection{显式方法与隐式方法}
\noindent\textit{依据：课件第46--52页；Anderson 第4.4节、第9.3节。}

\textbf{显式方法}由已知时间层直接算出新时间层；
\textbf{隐式方法}的新时间层出现在需要求解的代数关系中，通常须联立求解，非线性时还需迭代。
“一条方程有几个未知数”只是某些差分例子的直观表现，不是普遍定义。

对一维热方程 $T_t=\alpha T_{xx}$（常数 $\alpha>0$），定义扩散数
\[
 r=\frac{\alpha\Delta t}{\Delta x^2},\qquad
 \delta_{xx}T_i=\frac{T_{i+1}-2T_i+T_{i-1}}{\Delta x^2}.
\]
FTCS（时间向前、空间中心）为
\begin{equation}
 T_i^{n+1}=T_i^n+r(T_{i+1}^n-2T_i^n+T_{i-1}^n),
 \qquad O(\Delta t+\Delta x^2).
 \label{eq:cfd-heat-ftcs}
\end{equation}
Crank--Nicolson（CN）对新旧时间层的空间算子取平均：
\begin{equation}
 \frac{T_i^{n+1}-T_i^n}{\Delta t}
 =\frac\alpha2(\delta_{xx}T_i^{n+1}+\delta_{xx}T_i^n),
 \qquad O(\Delta t^2+\Delta x^2).
 \label{eq:cfd-cn}
\end{equation}
将未知量移至左端：
\begin{equation}
 -\frac r2T_{i-1}^{n+1}+(1+r)T_i^{n+1}-\frac r2T_{i+1}^{n+1}
 =\frac r2T_{i-1}^n+(1-r)T_i^n+\frac r2T_{i+1}^n.
 \label{eq:cfd-cn-tridiagonal}
\end{equation}
给定两端温度时，边界贡献并入右端；全部内部节点组成三对角系统，可用 Thomas 算法以 $O(N)$ 代价求解。
周期边界对应循环三对角系统，需采用相应修正。
显式格式单步简单，但通常受时间步限制；隐式格式可减轻某些稳定性限制，却增加线性／非线性求解代价。
\textbf{隐式不自动等于无条件稳定，无条件稳定也不等于任意大时间步都准确。}

\subsubsection{误差与稳定性分析}
\label{sec:cfd-stability}
\noindent\textit{依据：课件第53--72页；Anderson 第4.5节、第6.2--6.3节、第6.6节。}

\paragraph{A. 误差传播与 von Neumann 分析}
记连续问题的精确解在网格上的取值为 $u$，离散方程的精确解为 $D$，计算机结果为 $N$。
离散误差为 $D-u$，舍入误差为 $\epsilon=N-D$，总误差是两者之和。
若线性格式写为 $D^{n+1}=Q D^n+b^n$，则
\begin{equation}
 \epsilon^{n+1}=Q\epsilon^n+\eta^n,
 \label{eq:cfd-error-propagation}
\end{equation}
$\eta^n$ 是本步新引入的舍入／代数求解误差。
\textbf{忽略新的误差注入时}，扰动满足齐次差分方程；不能说实际每步产生的舍入误差严格满足无源方程。

稳定性的本质是在固定时间区间 $0\le n\Delta t\le T$ 内，
传播算子 $\|Q^n\|\le C_T$，且 $C_T$ 不随网格加密而发散。
误差不必逐步衰减；无耗散格式也可以稳定。
对\textbf{线性、常系数、均匀网格，且周期或无限域}的问题，可作 Fourier 展开：
\[
 \epsilon_i^n=\sum_m\widehat\epsilon_m^n e^{\mathrm i k_mx_i},\qquad
 \widehat\epsilon_m^{n+1}=g(\theta_m)\widehat\epsilon_m^n,
 \qquad\theta_m=k_m\Delta x.
\]
空间平移使各模态独立，因此可代入单个模式
$\epsilon_i^n=g^n e^{\mathrm i i\theta}$ 求放大因子，再检查\textbf{所有}可解析波数
$\theta\in[-\pi,\pi]$。不是所有模态有相同增长率，也不是只检查某一个选定波数。
对下述无物理增长的标量单步格式，常用 $|g(\theta)|\le1$ 的不放大判据；
更一般的有限时间稳定性允许 $|g|\le1+C\Delta t$。
多步格式还须检查全部根及单位圆重根，系统还须检查矩阵幂的一致有界性。
变系数、非线性和边界闭合可能引入额外不稳定，不能只靠这一分析下结论。

热方程 FTCS 代入后，
\[
 g=1+r(e^{\mathrm i\theta}-2+e^{-\mathrm i\theta})
   =1-4r\sin^2(\theta/2).
\]
故 $|g|\le1$ 对所有 $\theta$ 成立的条件是 $\boxed{0\le r\le1/2}$。
CN 则有
\begin{equation}
 g_{\mathrm{CN}}=\frac{1-2r\sin^2(\theta/2)}{1+2r\sin^2(\theta/2)},
 \qquad |g_{\mathrm{CN}}|\le1\quad(r\ge0).
 \label{eq:cfd-cn-amplification}
\end{equation}
这是上述热方程的无条件稳定结论。大 $r$ 时高频模态可能 $g<0$ 且接近 $-1$，
因而出现衰减很慢的交替振荡；稳定不保证单调、无振荡或准确。

\paragraph{B. 对流方程、FTCS 与 CFL}
以下取 $u_t+cu_x=0$，$c$ 为常数，定义带符号 Courant 数
\begin{equation}
 \nu=\frac{c\Delta t}{\Delta x}.
 \label{eq:cfd-courant}
\end{equation}
$|\nu|$ 是一个时间步内信息传播距离与网格间距之比。
这里用 $\nu$ 表示课件中的 CFL 数，避免与特征值 $\lambda$ 混淆；扩散数仍记为 $r$。
FTCS 为
\[
 u_i^{n+1}=u_i^n-\frac\nu2(u_{i+1}^n-u_{i-1}^n),\qquad
 g=1-\mathrm i\nu\sin\theta,
 \qquad |g|^2=1+\nu^2\sin^2\theta.
\]
任何非零 $\nu$ 都不能使所有模态满足 $|g|\le1$。
在通常固定非零 Courant 数、$\Delta t\propto\Delta x$ 的对流加密路径下，FTCS 不稳定，
课件称为“无条件不稳定”。严格按有限时间有界定义，若额外取
$\Delta t=O(\Delta x^2)$ 可限制累计放大；这不改变它缺少对流尺度稳定区间、通常不宜用于纯对流的结论。

\textbf{CFL 必要条件：数值依赖域必须包含连续问题的依赖域。}
三点显式格式每步最多跨一个网格，因此其信息传播上限为 $\Delta x/\Delta t$，
须有 $|c|\le\Delta x/\Delta t$，即 $|\nu|\le1$。
如图~\ref{fig:cfd-characteristics} 所示，对双向波动，物理依赖区间
$[x_0-ct_0,x_0+ct_0]$ 必须落在数值依赖区间内；对单向对流则检查初始脚点 $x_0-ct_0$。
\textbf{CFL 是必要条件，并非一般充分条件}：FTCS 即使满足 $|\nu|\le1$ 仍可能不稳定。
系统应采用最大特征速度，Euler 方程常涉及 $|v_n|+a_s$；
多维、多阶段及隐式格式的限制须按实际格式分析，不存在通用的“CFL 总要小于1”。

\paragraph{C. Leap-Frog（CTCS）格式}
时间和空间都用中心差分：
\begin{equation}
 u_i^{n+1}=u_i^{n-1}-\nu(u_{i+1}^n-u_{i-1}^n),
 \qquad O(\Delta t^2+\Delta x^2).
 \label{eq:cfd-leapfrog}
\end{equation}
代入 Fourier 模式，得到
\[
 g^2+2\mathrm i\nu\sin\theta\,g-1=0,\qquad
 g_\pm=-\mathrm i\nu\sin\theta\pm\sqrt{1-\nu^2\sin^2\theta}.
\]
当 $|\nu|<1$ 时，两根模长均为1：没有幅值耗散，但有相位误差。
需要两个时间层，必须用另一个一致格式计算第一步；除物理解支外，还有随步数交替的计算模态。
\textbf{端点需注意：}$|\nu|=1$ 时某些波数产生单位圆重根，任意两层扰动可出现
$n g^n$ 增长；教材常写 $|\nu|\le1$ 的模长条件，严格稳健计算通常取
$|\nu|\le\nu_{\max}<1$ 并采用相容启动。

同一个 Leap-Frog 用于热方程则有
$g^2+8r\sin^2(\theta/2)g-1=0$。
对 $r>0$ 的非零波数，必有一根模长大于1，故不稳定。
稳定性结论必须与\emph{所离散的方程}一起记忆。

\paragraph{D. Lax（Lax--Friedrichs）格式}
将对流 FTCS 右端的 $u_i^n$ 改为邻点平均：
\begin{equation}
 u_i^{n+1}=\frac{u_{i+1}^n+u_{i-1}^n}{2}
          -\frac\nu2(u_{i+1}^n-u_{i-1}^n).
 \label{eq:cfd-lax}
\end{equation}
其放大因子及模长为
\[
 g=\cos\theta-\mathrm i\nu\sin\theta,\qquad
 |g|^2=1-(1-\nu^2)\sin^2\theta.
\]
故稳定条件为 $|\nu|\le1$。它通过邻点平均引入数值耗散，会抹平陡峭梯度；
在固定非零 $\nu$ 下整体为一阶精度。
其一致性还要求 $\Delta x^2/\Delta t\to0$；固定空间网格时，盲目减小时间步反而会增加平均操作造成的扩散。

\paragraph{E. 用空间导数替换时间导数：Lax--Wendroff}
先在时间上作二阶 Taylor 展开，再用 PDE 消去时间导数：
\[
 u^{n+1}=u^n+\Delta t\,u_t+\tfrac12\Delta t^2u_{tt}+O(\Delta t^3),
 \qquad u_t=-cu_x,\quad u_{tt}=c^2u_{xx}.
\]
最后用空间中心差分，得到
\begin{equation}
 u_i^{n+1}=u_i^n-\frac\nu2(u_{i+1}^n-u_{i-1}^n)
 +\frac{\nu^2}{2}(u_{i+1}^n-2u_i^n+u_{i-1}^n).
 \label{eq:cfd-lax-wendroff}
\end{equation}
这是时间、空间二阶格式。
\begin{align*}
 g&=1-\mathrm i\nu\sin\theta+\nu^2(\cos\theta-1),\\
 |g|^2&=1-4\nu^2(1-\nu^2)\sin^4(\theta/2).
\end{align*}
因此稳定区间为 $\boxed{|\nu|\le1}$，不存在必须满足 $|\nu|>0.5$ 的普遍下限。
它比 Lax 更能保留光滑波形，但不是单调格式，在间断附近可能出现色散振荡。
对这个常速线性模型，$\nu=\pm1$ 时退化为恰好平移一个网格。

\paragraph{补充：两步 Lax--Wendroff 与 MacCormack}
对守恒律 $u_t+f(u)_x=0$，两步 Lax--Wendroff（Richtmyer）写为
\begin{align*}
 u_{i+1/2}^{n+1/2}
 &=\tfrac12(u_{i+1}^n+u_i^n)
   -\frac{\Delta t}{2\Delta x}[f(u_{i+1}^n)-f(u_i^n)],\\
 u_i^{n+1}
 &=u_i^n-\frac{\Delta t}{\Delta x}
       [f(u_{i+1/2}^{n+1/2})-f(u_{i-1/2}^{n+1/2})].
\end{align*}
MacCormack 则在原节点预估并反向校正：
\begin{align*}
 u_i^*&=u_i^n-\frac{\Delta t}{\Delta x}[f(u_{i+1}^n)-f(u_i^n)],\\
 u_i^{n+1}&=\tfrac12\left\{u_i^n+u_i^*
      -\frac{\Delta t}{\Delta x}[f(u_i^*)-f(u_{i-1}^*)]\right\}.
\end{align*}
二者对光滑解可达二阶；当 $f(u)=cu$ 时都化为式~\eqref{eq:cfd-lax-wendroff}。
对非线性系统不能直接照搬这个标量稳定性证明，间断附近通常还需适当的耗散或限制器。

\paragraph{F. 一致性、稳定性与收敛性}
\begin{itemize}
 \item \textbf{一致性（consistency）：}把足够光滑的 PDE 精确解代入离散算子，按微分方程量纲定义的截断误差随网格细化趋于0。
 \item \textbf{稳定性（stability）：}初值、舍入或强迫中的扰动在指定时间区间内受到与网格无关的控制。
 \item \textbf{收敛性（convergence）：}沿规定的 $\Delta t$、$\Delta x$ 加密关系，离散解在相应范数下趋于连续问题的解。
\end{itemize}
Lax 等价定理针对\textbf{适定的线性初值问题及其一致的线性差分逼近}，在相应函数空间与范数下给出
\[
 \boxed{\text{一致性}+\text{稳定性}\ \Longrightarrow\ \text{收敛性}},
 \qquad\text{且在这些前提下，稳定性与收敛性等价。}
\]
有边界时还必须把边界处理纳入稳定性讨论。
该定理不能直接推广到任意非线性 PDE；非线性守恒律还涉及离散守恒、紧性与熵条件等。

\paragraph{G. 修正方程、耗散与色散}
修正方程把差分格式展开为原 PDE 加上网格相关的高阶项，用于解释误差的作用。
它是光滑解下的渐近分析，有限项截断不能取代对所有网格波数的稳定性检验。
对纯对流，以固定 Courant 数进行展开，几个主导项为
\begin{align}
 \text{FTCS:}\quad u_t+cu_x
   &=-\frac{c^2\Delta t}{2}u_{xx}+\cdots,\\
 \text{迎风（$c>0$）:}\quad u_t+cu_x
   &=\frac{c\Delta x}{2}(1-\nu)u_{xx}+\cdots,\\
 \text{Lax:}\quad u_t+cu_x
   &=\frac{\Delta x^2}{2\Delta t}(1-\nu^2)u_{xx}+\cdots,\\
 \text{Lax--Wendroff:}\quad u_t+cu_x
   &=\frac{c\Delta x^2}{6}(\nu^2-1)u_{xxx}+\cdots.
 \label{eq:cfd-modified}
\end{align}
这里迎风格式为 $u_i^{n+1}=u_i^n-\nu(u_i^n-u_{i-1}^n)$，稳定范围为 $0\le\nu\le1$；
$c<0$ 时改用向前空间差分，其数值扩散系数是
$|c|\Delta x(1-|\nu|)/2$。
FTCS 的负扩散解释了误差增长；迎风与 Lax 的正扩散抹平梯度；
Lax--Wendroff 的主导三阶项改变波的相位。
\begin{itemize}
 \item \textbf{耗散（dissipation）：}不同波数的振幅被衰减。它可抑制高频噪声，也会使激波或接触间断变宽。
 \item \textbf{色散（dispersion）：}不同波数以不同的数值相速度传播，导致波形畸变与振荡。
       课件第68页把奇数阶导数的作用再次写成 dissipation，应改为 dispersion。
\end{itemize}
对实常系数的修正方程，偶数阶空间导数改变振幅，奇数阶空间导数改变相位；
偶数阶项究竟耗散还是反耗散，还取决于符号。具体取
\begin{equation}
 u_t+cu_x=\epsilon_2u_{xx}-\epsilon_4u_{xxxx},
 \qquad\epsilon_2,\epsilon_4\ge0.
 \label{eq:cfd-artificial-viscosity}
\end{equation}
因为 $\partial_{xx}e^{\mathrm ikx}=-k^2e^{\mathrm ikx}$、
$\partial_{xxxx}e^{\mathrm ikx}=k^4e^{\mathrm ikx}$，振幅的增长率是
$-\epsilon_2k^2-\epsilon_4k^4\le0$。
因此右端二阶项取正系数，四阶项取负系数才是耗散；四阶耗散对短波更有选择性。
不能用“二阶导数总小于0”解释稳定性，它的点值并不总为负。
四阶导数的常用中心近似为
\[
 u_{xxxx}(x_i)\approx
 \frac{u_{i+2}-4u_{i+1}+6u_i-4u_{i-1}+u_{i-2}}{\Delta x^4}.
\]
人工黏性是人为或由离散隐含引入的数值耗散，区别于流体的物理黏性。
加入人工项后仍须检验完整格式的稳定性；显式推进还会受
$\epsilon_2\Delta t/\Delta x^2$、$\epsilon_4\Delta t/\Delta x^4$ 的限制。
若目标仍是原无黏方程，人工项必须在加密极限中以适当方式消失，并宜以通量差形式加入以保持离散守恒。

\paragraph{补充：对流扩散与算子分裂}
对 $u_t+cu_x=\alpha u_{xx}$，未分裂 FTCS 的放大因子为
$g=1-4r\sin^2(\theta/2)-\mathrm i\nu\sin\theta$，稳定条件是
\[
 0\le r\le\tfrac12,\qquad \nu^2\le2r.
\]
这说明对流与扩散的限制应一起检查，不能简单套用两个独立结论。
若写成 $u_t=Lu+Nu$，Lie 分裂依次求解 $u_t=Lu$ 和 $u_t=Nu$ 各推进 $\Delta t$，
组合近似一个完整步，通常时间一阶；Strang 分裂按
$L(\Delta t/2)\to N(\Delta t)\to L(\Delta t/2)$ 对称组合，
在解与算子足够光滑、子问题求解足够准确时可达时间二阶。
子问题的右端符号必须与原式一致；分裂不会自动消除子步骤的稳定性与边界误差。

\clearpage
\subsubsection{计算网格}
\label{sec:cfd-grid}
\noindent\textit{依据：课件第37--39、73--83页；Anderson 第5.2--5.7节。}

\paragraph{网格类型与质量}
结构网格具有规则的逻辑索引，可构成 O、C、H 型贴体网格；非结构网格由显式的单元连接关系组织，
适合复杂几何与局部加密。网格应准确表示边界，在边界层、激波等高梯度区域加密，
使尺寸变化平缓并控制倾斜程度。课件的网格夹角 $60^\circ$--$120^\circ$ 是经验范围，
不是普适稳定性定理；应结合离散方法检查正交性、扭曲度与网格收敛性。

\paragraph{物理平面到计算平面}
以\textbf{固定不随时间运动}的二维光滑贴体网格为例，建立
\[
 (x,y)\longleftrightarrow(\xi,\eta),\qquad
 x=x(\xi,\eta),\quad y=y(\xi,\eta).
\]
物理平面中的曲线、非均匀网格在计算平面中对应规则索引网格；
场量仍代表相同物理点，几何复杂性转移到度量系数中。
一阶导数的链式法则是
\begin{align*}
 \partial_\xi&=x_\xi\partial_x+y_\xi\partial_y,
 &\partial_\eta&=x_\eta\partial_x+y_\eta\partial_y,\\
 \partial_x&=\xi_x\partial_\xi+\eta_x\partial_\eta,
 &\partial_y&=\xi_y\partial_\xi+\eta_y\partial_\eta.
\end{align*}
定义两个互为倒数的行列式：
\begin{equation}
 D=\det\frac{\partial(x,y)}{\partial(\xi,\eta)}
   =x_\xi y_\eta-x_\eta y_\xi,
 \qquad
 J=\det\frac{\partial(\xi,\eta)}{\partial(x,y)}=\frac1D.
 \label{eq:cfd-jacobian}
\end{equation}
\textbf{本讲采用课件的 $J$ 约定；Anderson 第5章的 Jacobian 对应这里的 $D$。}
局部可逆要求 $D\ne0$，度量有限；实际网格还须无全局折叠。
以下取保向映射 $D>0$，于是 $\dd x\dd y=D\dd\xi\dd\eta=J^{-1}\dd\xi\dd\eta$。
仅要求“Jacobian 有限”不足以排除退化。

逆矩阵给出
\begin{equation}
 \begin{bmatrix}\xi_x&\xi_y\\\eta_x&\eta_y\end{bmatrix}
 =\begin{bmatrix}x_\xi&x_\eta\\y_\xi&y_\eta\end{bmatrix}^{-1}
 =J\begin{bmatrix}y_\eta&-x_\eta\\-y_\xi&x_\xi\end{bmatrix}.
 \label{eq:cfd-metrics}
\end{equation}
故
\begin{equation}
 \partial_x=J(y_\eta\partial_\xi-y_\xi\partial_\eta),\qquad
 \partial_y=J(-x_\eta\partial_\xi+x_\xi\partial_\eta).
 \label{eq:cfd-transformed-derivatives}
\end{equation}
二阶导数还要对度量求导，不能把变系数当常数直接平方。

\paragraph{曲线坐标下的二维 Euler 方程}
无体源时，物理平面中的方程为
\[
 \mathbf U_t+\mathbf F_x+\mathbf G_y=0,
 \quad
 \mathbf U=\begin{bmatrix}\rho\\\rho u\\\rho v\\E_t\end{bmatrix},\quad
 \mathbf F=\begin{bmatrix}\rho u\\\rho u^2+p\\\rho uv\\u(E_t+p)\end{bmatrix},\quad
 \mathbf G=\begin{bmatrix}\rho v\\\rho uv\\\rho v^2+p\\v(E_t+p)\end{bmatrix},
\]
其中 $E_t=\rho[e+(u^2+v^2)/2]$。代入式~\eqref{eq:cfd-transformed-derivatives} 并除以 $J$，得到
\[
 \frac1J\mathbf U_t
 +y_\eta\mathbf F_\xi-y_\xi\mathbf F_\eta
 -x_\eta\mathbf G_\xi+x_\xi\mathbf G_\eta=0.
\]
由于映射不随时间变，且混合偏导可交换，满足度量恒等式
\[
 \partial_\xi y_\eta-\partial_\eta y_\xi=0,\qquad
 -\partial_\xi x_\eta+\partial_\eta x_\xi=0.
\]
利用乘积法则，度量导数项相消，从而恢复强守恒形式：
\begin{equation}
 \boxed{\frac{\partial\widetilde{\mathbf U}}{\partial t}
       +\frac{\partial\widetilde{\mathbf F}}{\partial\xi}
       +\frac{\partial\widetilde{\mathbf G}}{\partial\eta}=0},
 \label{eq:cfd-curvilinear-euler}
\end{equation}
\begin{align}
 \widetilde{\mathbf U}&=\frac{\mathbf U}{J}=D\mathbf U,\\
 \widetilde{\mathbf F}&=\frac{\xi_x\mathbf F+\xi_y\mathbf G}{J}
                   =y_\eta\mathbf F-x_\eta\mathbf G,\\
 \widetilde{\mathbf G}&=\frac{\eta_x\mathbf F+\eta_y\mathbf G}{J}
                   =-y_\xi\mathbf F+x_\xi\mathbf G.
 \label{eq:cfd-transformed-fluxes}
\end{align}
定义逆变速度 $V^\xi=\xi_xu+\xi_yv$、$V^\eta=\eta_xu+\eta_yv$，则展开的通量为
\begin{equation}
 \widetilde{\mathbf F}=\frac1J
 \begin{bmatrix}
 \rho V^\xi\\\rho uV^\xi+\xi_xp\\\rho vV^\xi+\xi_yp\\(E_t+p)V^\xi
 \end{bmatrix},\qquad
 \widetilde{\mathbf G}=\frac1J
 \begin{bmatrix}
 \rho V^\eta\\\rho uV^\eta+\eta_xp\\\rho vV^\eta+\eta_yp\\(E_t+p)V^\eta
 \end{bmatrix}.
 \label{eq:cfd-expanded-curvilinear-fluxes}
\end{equation}
$V^\xi,V^\eta$ 含坐标尺度，通常不是沿单位法向的物理速度。
若有体源，右端相应为 $\mathbf S/J$。
例如 $x=2\xi,y=3\eta$ 时，$D=6,J=1/6$，应得到
$\widetilde{\mathbf U}=6\mathbf U$、$\widetilde{\mathbf F}=3\mathbf F$、
$\widetilde{\mathbf G}=2\mathbf G$，可用于自检是否把 $J$ 与 $1/J$ 混淆。
离散度量应兼容上述恒等式，以保持均匀来流；运动网格还需网格速度通量与几何守恒律，不能直接套用固定网格公式。

\paragraph{结构网格生成方法}
\begin{itemize}
 \item \textbf{代数／几何法：}由边界插值并用拉伸函数控制加密，生成快，但复杂边界下的正交性不易保证。
 \item \textbf{复变量法：}二维共形映射满足 Cauchy--Riemann 关系
       $\xi_x=\eta_y$、$\xi_y=-\eta_x$；导数非零处保角，适于构造正交网格。
       单独满足两个 Laplace 方程并不自动保证共形，不能直接推广为三维保角方法。
 \item \textbf{PDE 法：}求解椭圆型网格方程，如
       $\nabla^2\xi=P(\xi,\eta)$、$\nabla^2\eta=Q(\xi,\eta)$，
       由控制函数调节疏密并获得平滑网格；计算较贵，仍须检查网格是否折叠。
\end{itemize}

\paragraph{与数据驱动／降阶方法的联系（补充）}
空间离散把 PDE 变成高维动力系统 $\dot{\mathbf a}=\mathbf R(\mathbf a)$。
令 $\bar{\mathbf a}$ 为样本均值，收集中心化快照矩阵
$X=[\mathbf a(t_1)-\bar{\mathbf a},\ldots,\mathbf a(t_m)-\bar{\mathbf a}]$，
通过 $X=\Phi\Sigma\Psi^{\mathsf T}$ 的截断 SVD 提取 POD 基，
近似 $\mathbf a\approx\bar{\mathbf a}+\Phi_r\mathbf z$，再作 Galerkin 投影可得降阶模型。
若使用非均匀网格，物理积分内积应包含网格体积／Jacobian 权重。
低秩表示误差小并不自动保证降阶动力学稳定或守恒，仍须检验。
这里仅使用 Brunton 与 Kutz 第1章、第11.1节支撑 SVD、POD 和投影的联系；
不把这些内容作为前述 CFD 稳定性结论的依据。\cite{cfd:brunton}

\begin{thebibliography}{9}
 \bibitem{cfd:ppt}
 王兴建：《第一讲 CFD基础（第1、2周）-更新版》，2026年秋课程课件，83页。
 本讲主要依据第18--34页（方程与PDE）、第35--72页（离散与稳定性）、第73--83页（网格）。
 引用页码为PDF文件页序；对简化表述及明显笔误已在正文更正。

 \bibitem{cfd:anderson}
 John D. Anderson, Jr.：《计算流体力学基础及其应用》，所提供中文扫描版。
 重点参照第2章（控制方程）、第3章（PDE性质）、第4章（离散、显隐式与稳定性）、
 第5章（网格和坐标变换）、第6.2--6.3节（Lax--Wendroff与MacCormack）、
 第6.6节（耗散、色散及人工黏性）、第9.3节（Crank--Nicolson）。
 对应PDF文件页序约为37--77、79--96、98--126、129--161、163--168、173--179、303--313。

 \bibitem{cfd:brunton}
 Steven L. Brunton and J. Nathan Kutz:
 \textit{Data-Driven Science and Engineering: Machine Learning, Dynamical Systems, and Control}.
 Cambridge University Press, 2019.
 第1章、第11.1节；后者为所提供PDF第398页起，印刷页码375页起。
 仅用于本讲末尾关于SVD、POD及Galerkin降阶的补充。
\end{thebibliography}
